% !TEX root = ../../../practices.tex
% ==============================================================================
% EJERCICIO 1 - PRÁCTICA 0: CONTEO Y PROBABILIDAD
% Ubicación: practices/practice_00_counting_probability/exercises/ex_01.tex
% ==============================================================================
\IfFileExists{practices.tex}{%
    \documentclass[practices.tex]{subfiles}%
}{%
    \documentclass[../../../practices.tex]{subfiles}%
}

\begin{document}

\h2{Ejercicio 1}

Hay 20 candidatos para ocupar 3 cargos distintos. ¿De cuántas maneras diferentes se podrían ocupar los tres cargos?

\h3{Desarrollo y respuesta}
Para resolver este problema representamos los tres cargos como casillas que asignamos en orden sucesivo.

Como los cargos tienen funciones distintas, el orden de asignación es determinante. No es lo mismo que un candidato sea elegido para el primer puesto a que resulte electo para el segundo. Además, una misma persona no puede ocupar dos cargos a la vez, por lo que el proceso se efectúa sin repetición.

\begin{enumerate}
    \item Para el primer cargo disponemos de cualquiera de los 20 candidatos.
    \item Para el segundo cargo, al haber seleccionado ya a una persona, restan 19 postulantes disponibles.
    \item Para el tercer cargo, cubiertos los puestos previos, quedan 18 opciones posibles.
\end{enumerate}

El número total de asignaciones se obtiene aplicando el principio fundamental de la multiplicación mediante variaciones sin repetición de 20 elementos tomados de 3 en 3,
\begin{equation*}
    N = 20 \times 19 \times 18 = \frac{20!}{(20 - 3)!} = 6840.
\end{equation*}

\conclusion{Existen $6840$ maneras diferentes de ocupar los tres cargos distintos a partir del grupo de 20 candidatos.}

\end{document}
